\section*{الفصل \ref{ch:g11:nucleus-radioactivity} --- النواة والنشاط الإشعاعي}

\begin{solution}{exo:g11:nucleus-radioactivity:1}
${}^{16}_{8}\mathrm{O}$: $8$ p، $8$ n؛ ${}^{27}_{13}\mathrm{Al}$: $13$ p،
$14$ n؛ ${}^{60}_{27}\mathrm{Co}$: $27$ p، $33$ n؛
${}^{235}_{92}\mathrm{U}$: $92$ p، $143$ n؛ ${}^{3}_{1}\mathrm{H}$: $1$ p،
$2$ n. (وفي كل مرة $N = A - Z$.)
\end{solution}

\begin{solution}{exo:g11:nucleus-radioactivity:2}
الكربونات الثلاثة ($Z = 6$) \omterm{def:g11:nucleus-radioactivity:notation}{نظائر} بعضها لبعض.
أما ${}^{14}_{6}\mathrm{C}$ و${}^{14}_{7}\mathrm{N}$ فيشتركان في $A = 14$ لا في
$Z$: فهما عنصران مختلفان (بسحابتَي إلكترونات مختلفتين)، ومن ثم ليسا نظيرين.
\end{solution}

\begin{solution}{exo:g11:nucleus-radioactivity:3}
$A = 210 - 4 = 206$، $Z = 84 - 2 = 82$:
${}^{210}_{84}\mathrm{Po} \to {}^{206}_{82}\mathrm{Pb} +
{}^{4}_{2}\mathrm{He}$ --- فالوليدة هي الرصاص-206.
\end{solution}

\begin{solution}{exo:g11:nucleus-radioactivity:4}
${}^{60}_{27}\mathrm{Co} \to {}^{60}_{28}\mathrm{Ni} +
{}^{\;0}_{-1}\mathrm{e}$ و
${}^{131}_{53}\mathrm{I} \to {}^{131}_{54}\mathrm{Xe} +
{}^{\;0}_{-1}\mathrm{e}$: إذ إنّ $A$ لا يتغير، و$Z$ يزيد واحدًا.
\end{solution}

\begin{solution}{exo:g11:nucleus-radioactivity:5}
$16$ يومًا $= 2\,T_{1/2}$: $800/4 = \qty{200}{MBq}$. و$32$ يومًا
$= 4\,T_{1/2}$: $800/16 = \qty{50}{MBq}$.
\end{solution}

\begin{solution}{exo:g11:nucleus-radioactivity:6}
(أ) $\alpha$ (ثقيل، ومن ثم لا يكاد يُزاح)؛ (ب) $\beta^-$؛ (ج) $\gamma$
(متعادل نافذ)؛ (د) $\beta^+$ (\omterm{rem:g11:nucleus-radioactivity:betaplus}{البوزيترون}: بخفّة (ب) نفسها، وبشحنة
معاكسة).
\end{solution}

\begin{solution}{exo:g11:nucleus-radioactivity:7}
${}^{222}_{86}\mathrm{Rn} \to {}^{218}_{84}\mathrm{Po} +
{}^{4}_{2}\mathrm{He}$؛ \
${}^{218}_{84}\mathrm{Po} \to {}^{214}_{82}\mathrm{Pb} +
{}^{4}_{2}\mathrm{He}$؛ \
${}^{214}_{82}\mathrm{Pb} \to {}^{214}_{83}\mathrm{Bi} +
{}^{\;0}_{-1}\mathrm{e}$.
\end{solution}

\begin{solution}{exo:g11:nucleus-radioactivity:8}
للفلور-18 $9$ نيوترونًا حيث للفلور-19 المستقر $10$:
فهو فقير بالنيوترونات، تحت الشريط، ومن ثم $\beta^+$:
${}^{18}_{9}\mathrm{F} \to {}^{18}_{8}\mathrm{O} + {}^{0}_{+1}\mathrm{e}$.
\end{solution}

\begin{solution}{exo:g11:nucleus-radioactivity:9}
$24$ ساعة $= 4\,T_{1/2}$: $500/16 \approx \qty{31}{MBq}$. وبما أن
$2^9 = 512$، تُنزل تسعة أعمار نصف النشاطَ إلى $500/512 < \qty{1}{MBq}$: أي بعد
$9 \times 6 = \qty{54}{h}$.
\end{solution}

\begin{solution}{exo:g11:nucleus-radioactivity:10}
$\frac18 = \frac1{2^3}$: أي ثلاثة أعمار نصف، $3 \times \num{5700} =
\num{17100}$ سنة. أما عظم ديناصور فعمره نحو $\num{17500}$ عمر نصف:
فيبقى منه جزء قدره $1/2^{17500}$ --- أي ولا \omterm{def:g11:fundamental-interactions:ladder}{ذرة} واحدة من الكربون-14 (فغرام من
الكربون لا يحوي إلا نحو $\num{5e22}$ \omterm{def:g11:fundamental-interactions:ladder}{ذرة}).
\end{solution}

\begin{solution}{exo:g11:nucleus-radioactivity:11}
${}^{14}_{6}\mathrm{C}$: $8$ نيوترونًا مقابل $6$--$7$ في الكربون المستقر،
أي فوق الشريط، ومن ثم $\beta^-$. و${}^{11}_{6}\mathrm{C}$: $5$ نيوترونًا،
فقير بالنيوترونات، ومن ثم $\beta^+$. و${}^{238}_{92}\mathrm{U}$: $Z = 92 > 83$، ثقيل
جدًّا، ومن ثم $\alpha$.
\end{solution}

\begin{solution}{exo:g11:nucleus-radioactivity:12}
لا يغيّر $A$ إلا $\alpha$: أي $238 - 206 = 32 = 4x$، ومنه $x = 8$ جسيم \omterm{def:g11:nucleus-radioactivity:abg}{ألفا}. وهي
تنزع $16$ من $Z$: فيعطي $92 - 16 + y = 82$ عدد $y = 6$ تفكك بيتا-ناقص.
\end{solution}

\begin{solution}{exo:g11:nucleus-radioactivity:13}
$6.4/0.1 = 64 = 2^{6}$: أي ستة أعمار نصف، $6 \times 5.3 = 31.8 \approx 32$
سنة في تخزين مدرَّع.
\end{solution}

\begin{solution}{exo:g11:nucleus-radioactivity:14}
في اليوم: $8000 \times 86400 \approx \num{6.9e8}$ تفكك. وعلى مدى $80$ سنة:
$\num{6.9e8} \times 365 \times 80 \approx \num{2e13}$. وقد جرت الحياة دائمًا
على هذه الخلفية --- \omterm{def:g11:nucleus-radioactivity:activity}{فالبيكريل} الواحد \omterm{def:g11:fundamental-interactions:ladder}{ذرة} واحدة في \omterm{def:g10:orders-of-magnitude:unit}{الثانية} من بين
$\sim 10^{27}$ في الجسم؛ \omterm{def:g10:orders-of-magnitude:unit}{فالوحدة} ضئيلة، وعدد البيكريلات المخيف
الوقع قد يساوي بضع موزات.
\end{solution}

\begin{solution}{exo:g11:nucleus-radioactivity:15}
$\num{25000} \times 86400 \approx \num{2.2e9}$ تفكك في اليوم. والعشر سنوات
$10/432 \approx 2\%$ من \omterm{def:g11:nucleus-radioactivity:halflife}{عمر النصف}: أي أن ما تفكك أقلّ بكثير من النصف، ومن ثم
تفرض الإلكترونيات والغبار، لا المنبع، الاستبدال. وفي الخارج، تموت أشعة
$\alpha$ في هواء الحجرة وفي الغلاف؛ أما الغبار المستنشَق فيركن
المُصدِر في الرئتين، حيث يفرغ كل $\alpha$ طاقته في نسيج
حيّ --- أي جرعة كبيرة (\omterm{rem:g11:nucleus-radioactivity:dose}{بالسيفرت}) من \omterm{def:g11:nucleus-radioactivity:activity}{نشاط} متواضع.
\end{solution}

\begin{solution}{pb:g11:nucleus-radioactivity:1}
\textbf{1.} ${}^{12}_{6}\mathrm{C}$: $6$ p، $6$ n؛ و${}^{14}_{6}\mathrm{C}$:
$6$ p، $8$ n. ولا ترى الكيمياء إلا سحابة الإلكترونات، التي يثبّتها $Z = 6$:
فالسلوك متطابق.

\textbf{2.} الأكل والتنفس يجدّدان باستمرار كربون الجسم الحيّ
عند النسبة الجوية؛ وعند الموت يتوقف الأخذ ويجري التفكك
بلا معارض.

\textbf{3.} ${}^{14}_{6}\mathrm{C} \to {}^{14}_{7}\mathrm{N} +
{}^{\;0}_{-1}\mathrm{e}$ ($\beta^-$). وأشعة $\beta$ هذه ليّنة:
فأمتار من الكتّان الملفوف توقفها، فلا يفلت من المومياء إلا القليل --- ولهذا
تُقاس عيّنة كربون صغيرة تُؤخذ من الشيء نفسه.

\textbf{4.} $6.8/13.6 = \frac12$: أي عمر نصف واحد. فعمر المومياء نحو
\num{5700} سنة.

\textbf{5.} $13.2/13.6 \approx 0.97$: فالكتّان حديث أساسًا ---
وعمره بضعة قرون على الأكثر. أي تزوير.

\textbf{6.} $13.6/2^{10} = 13.6/1024 \approx 0.013$ تفكك في الدقيقة ---
أي عدّة واحدة كل $75$ دقيقة لكل غرام، غارقة في الخلفية الطبيعية. فعشرة
أعمار نصف، أي نحو \num{57000} سنة، هي الأفق العملي.

\textbf{7.} $92$ بروتونًا، و$146$ نيوترونًا. ويؤثر تنافر كولوم بين
كل أزواج البروتونات عبر \omterm{def:g11:fundamental-interactions:ladder}{النواة} بينما لا يربط الغراء القوي إلا
الجيران: فلا تصمد النوى الثقيلة إلا بفائض من النيوترونات، أي غراء بلا
تنافر.

\textbf{8.} ${}^{238}_{92}\mathrm{U} \to {}^{234}_{90}\mathrm{Th} +
{}^{4}_{2}\mathrm{He}$.

\textbf{9.} ${}^{234}_{90}\mathrm{Th} \to {}^{234}_{91}\mathrm{Pa} +
{}^{\;0}_{-1}\mathrm{e}$، ثم ${}^{234}_{91}\mathrm{Pa} \to
{}^{234}_{92}\mathrm{U} + {}^{\;0}_{-1}\mathrm{e}$: فيعاود اليورانيوم الظهور،
يورانيومًا-234.

\textbf{10.} $A$: $238 - 206 = 32$، ومنه $8$ خطوة $\alpha$؛ و$Z$:
$92 - 16 + y = 82$، ومنه $6$ خطوة $\beta^-$.

\textbf{11.} $\num{4.5e9}$ سنة عمر نصف واحد: فيبقى $\frac12$. وخمسة
مليارات سنة أخرى $\approx$ عمر نصف إضافي واحد: أي نحو $\frac14$.

\textbf{12.} تقود البراكين وتكتونية الصفائح (القارات الجارفة،
والزلازل) والحرارة الجوفية: أي جيولوجيا سطح كوكب يدفّئه
باطنه.

\textbf{13.} ${}^{241}_{95}\mathrm{Am} \to {}^{237}_{93}\mathrm{Np} +
{}^{4}_{2}\mathrm{He}$.

\textbf{14.} $\num{3.3e4}$ تفكك في \omterm{def:g10:orders-of-magnitude:unit}{الثانية}؛
و$\num{33000} \times 86400 \approx \num{2.9e9}$ في اليوم.

\textbf{15.} كل $\alpha$ يؤيّن الهواء الذي يعبره؛ وتحمل الأيونات
تيارًا ضئيلًا بين مسريَي الحجرة. وتلتقط جسيمات الدخان
الأيونات، فيهبط التيار، وينطلق الإنذار.

\textbf{16.} تنفق أشعة $\alpha$ كل طاقتها داخل الحجرة ولا
يفلت منها شيء من الغلاف: أي تأيين أقصى وتسرّب صفري. أما منبع
$\gamma$ \omterm{def:g11:nucleus-radioactivity:activity}{بالنشاط} نفسه فلا يكاد يؤيّن هواء الحجرة البتة، ويشعّ
الغرفة كلها بدل ذلك.

\textbf{17.} العشر سنوات $10/432 \approx 2\%$ من \omterm{def:g11:nucleus-radioactivity:halflife}{عمر النصف}: \omterm{def:g11:nucleus-radioactivity:activity}{فالنشاط}
لم يتغير أساسًا. والحلقات الأضعف هي الغبار والحشرات و
الإلكترونيات --- \omterm{def:g11:fundamental-interactions:ladder}{فالنواة} تعمّر أطول من الجهاز.

\textbf{18.} $\num{9e9}$ سنة $= 2$ عمر نصف: فيبقى $\frac14$.
ويورانيوم الأرض ما يزال وفيرًا، ومن ثم صُنع قبل بضعة
أعمار نصف على الأكثر --- فقد تكثّف الكوكب من رماد نجمي طازج، بعد
مليارات السنين من أولى النجوم.

\textbf{19.} ميت جيولوجيًّا: بباطن بارد، ولا بركنة ولا تكتونية
صفائح تعيد تدوير الهواء والصخر، ولا \omterm{def:g11:fundamental-interactions:ladder}{نواة} تخضّ --- ومن ثم لا
درع مغناطيسي في وجه الريح الشمسية.

\textbf{20.} كل \omterm{def:g11:fundamental-interactions:ladder}{ذرة} فيك أثقل من الهيليوم صُنعت في نجم محتضر،
وبقايا ذلك الأتون غير المستقرة ما تزال تدفّئ الأرض التي تقف
عليها: غبار نجوم، يبقيه عمر نصفه حيًّا.
\end{solution}
